% Standalone section -- paste into the weekly report, or \input it.
% A high-level account of coronary hemodynamics, written to explain the week-2 scope rather than to
% teach fluid mechanics: why flow is the reason to measure shape, and how that reason bears on each
% choice made this week (LM/LAD/LCX, tortuosity, bifurcation angle with diameters, a healthy
% population, topology repair). The full clinical account (the WSS equation, FFR, CT-FFR) is in
% thesis/clinical_background/05_hemodynamics.tex and is not repeated here.
%
% Citation keys resolve against thesis/refs.bib. griffo2026wss, zhang2024curvature and
% taylor2024murray were not cited by any .tex file before this one, so refs_clean.bib has to be
% regenerated after editing: python scripts/make_refs_clean.py
%
% CITATION WANTED: the bend-relative-to-width point is the Dean number (Dean 1927). There is no entry
% in refs.bib yet, so it is stated without a citation below; add one before this goes into the thesis.
% SECOND-HAND: the helical-flow statement and the split simulation literature on bifurcation angle
% come from shen2026geometry's Table 7, not from the primary papers. Cite for framing only.

\section*{Hemodynamics}

\textbf{Why flow is the reason to measure shape.}
Blood flowing through an artery drags on its lining, and that drag -- the wall shear stress -- is
not the same everywhere.
Every part of the coronary tree sees the same cholesterol and the same blood pressure, yet plaque
starts in particular places: where the drag is low, or reverses direction over the heartbeat, the
lining turns inflamed and plaque-prone \cite{chatzizisis2007ess}.
Those places are set by the shape of the vessel (\Cref{fig:week2-shape-factors}).
Tracing the flow through human coronary trees after death put almost every plaque on the outer wall
of a vessel leaving a bifurcation or on the inner wall of a bend \cite{asakura1990flow}, and the
drag also depends steeply on width: at the same flow, it scales with one over the radius cubed.
Shape sets the flow, and the flow sets where disease begins.

\begin{figure}[htbp]
  \centering
  \includegraphics[width=0.9\textwidth]{content/background/figures/shape_factors_shen2026.png}
  \caption{The shape features this note argues for, and the flow they act on
    \cite{shen2026geometry}. (a) A left coronary tree with its main bifurcation and two tortuous
    segments -- the two places low shear is expected. (b) Bifurcation angles, inflow angle and
    branch diameters. (c) Curvature $k = 1/R$ is local, while the tortuosity index $\tau = L/D$ is
    one number per segment. (d) Simulated near-wall shear, velocity, secondary and helical flow.
    Figure from Shen et al., Arch.\ Comput.\ Methods Eng.\ 2026,
    doi:10.1007/s11831-026-10530-w, CC BY 4.0.}
  \label{fig:week2-shape-factors}
\end{figure}

\textbf{Why I measure shape rather than simulate flow.}
Computing shear directly means a flow simulation per patient, with boundary conditions a CCTA scan
cannot supply: even a study of normal arteries had to give every vessel the same inlet pressure and
outlet flow \cite{schultz2023ess}.
That cost is why the largest whole-tree study of geometry and flow I have found covers 39 left
coronary trees \cite{zhang2024curvature}.
It is also increasingly avoidable. A network given only the vessel geometry reproduces the simulated
shear pattern well once each vessel is normalized by its own average, the correlation with the
simulation rising from 0.67 to 0.89 \cite{griffo2026wss}: the relative pattern is carried by the
shape, the absolute level much less so.
So this thesis does not run flow simulations. Hemodynamics is the argument for \emph{which}
geometric features to measure and in what form -- local rather than one number per patient, relative
rather than absolute -- and the geometry is what I actually measure.

\textbf{What this means for the scope chosen this week.}
\begin{itemize}
  \item \textbf{LM, LAD and LCX.} The left main bifurcation is where the flow first divides, and the
    walls just downstream of it are one of the two places low shear is expected. It is also where the
    flow literature has looked most, so the evidence linking geometry to shear is densest there
    \cite{kashyap2022tortuosity, zhang2024curvature}. The field's own review points out the flip
    side: vulnerable plaque is mostly found outside the left main region \cite{shen2026geometry},
    which is why these three vessels are a start rather than the whole scope.

  \item \textbf{Tortuosity.} A bend is the other low-shear site, and here hemodynamics chooses
    between metrics. In 127 patients without CAD, average absolute curvature was the tortuosity
    measure that best tracked the area of low shear across the left main branches, while the
    tortuosity index -- arc length over endpoint distance -- showed no relation at all ($p=0.86$)
    \cite{kashyap2022tortuosity}. In the 39-tree study, average curvature also separated
    stenosis-free from stenosed trees \cite{zhang2024curvature}. That is why my first implementation
    computes curvature, turning angle per mm, inflections, torsion and non-planarity alongside the
    index rather than the index alone.

    Three further points come from the flow. How much a bend disturbs the flow depends on how sharp
    it is \emph{relative to the vessel's width} (the Dean number, in fluid mechanics), so curvature
    has to be read against the lumen radius -- which is why the script measures radius as well.
    Out-of-plane twist drives helical flow, which several simulation studies propose as protective
    rather than harmful \cite{shen2026geometry}, so torsion should be kept as its own number rather
    than folded into a single tortuosity score. And flow acts locally: across 1010 patients, those
    with tortuous arteries were if anything \emph{less} likely to have CAD \cite{li2011tortuosity},
    whereas segment-level studies find more tortuosity where the plaque sits \cite{shen2026geometry},
    so tortuosity should be measured per segment, not once per patient.

    The catch is that curvature and torsion are derivatives, so their values depend on how much the
    centerline is smoothed before they are taken, while the index barely moves. The metrics
    hemodynamics favors are the ones hardest to make reproducible, and that trade-off is where I
    expect most of the method work to be.

  \item \textbf{Bifurcation angle, and why it is not enough on its own.} The angle changes where the
    flow separates from the wall, but how much blood each daughter vessel takes is set largely by
    its diameter, as Murray's law describes \cite{taylor2024murray}, and shear depends on that flow
    over the radius cubed. Two bifurcations with the same angle but different diameter ratios
    therefore load their walls differently -- the hemodynamic reason for wanting to extend the angle
    with the branch diameters. It may also be part of why simulation studies disagree on whether the
    angle matters at all \cite{shen2026geometry}.

  \item \textbf{A healthy population.} Shear acts before plaque exists, so the geometry that matters
    is the one before disease has remodeled it. The authors of the 39-tree study name exactly this
    gap: the role of flow in plaque \emph{onset} is underexplored, largely because longitudinal
    follow-up starting from a healthy state is lacking \cite{zhang2024curvature}. Repeat scans in
    the CGPS would be that follow-up.

  \item \textbf{Topology repair.} Flow needs a continuous path. A break in the predicted tree leaves
    part of a vessel with no path back to the ostium, and every feature measured along a path --
    tortuosity from ostium to tip, which branch comes off where -- is wrong for that vessel.
\end{itemize}

\textbf{What this does not claim.}
None of these features is a measurement of shear, and a feature that tracks low shear is a candidate
risk marker, not a predictor. Even simulated shear only moderately separated the lesions that went
on to cause a heart attack from those that did not, in patients who all had one (AUC 0.63--0.68)
\cite{griffo2026wss}, so shape features should not be expected to do better.
